% Responsable: Tahiny · Ronda 1 (sábado, hasta 12:00)

\subsection{Modelo en cascada}
\label{sec:modelo-cascada}

El modelo en cascada organiza el trabajo en fases sucesivas, con productos
revisados y aprobados antes de avanzar a la siguiente etapa. En ingeniería
de requisitos, su característica principal es establecer una especificación
suficientemente completa y validada antes de usarla como base del diseño y
la implementación. Esto facilita planificar entregables y conservar la
trazabilidad, pero también puede hacer costoso descubrir requisitos omitidos
en etapas tardías \cite{swebok2024,pohl2025}.

\subsection{Fases principales}
\begin{enumerate}
  \item \textbf{Delimitación y viabilidad.} Identificar el problema, los
  objetivos, las partes interesadas, el alcance, las restricciones y los
  riesgos iniciales. \emph{Resultado:} alcance y plan de trabajo acordados.

  \item \textbf{Obtención y análisis de requisitos.} Recopilar las
  necesidades de los actores, revisar el funcionamiento del entorno,
  resolver contradicciones y distinguir requisitos funcionales, de calidad
  y restricciones. \emph{Resultado:} requisitos analizados y priorizados.

  \item \textbf{Especificación y validación.} Documentar los requisitos
  con identificadores y criterios de aceptación; comprobar su coherencia,
  factibilidad, verificabilidad y correspondencia con las necesidades del
  cliente. Una vez aprobados, forman una \emph{línea base} sometida a
  control de cambios. \emph{Resultado:} especificación validada
  \cite{pohl2025,ireb2026}.

  \item \textbf{Diseño de la solución.} Definir arquitectura, módulos,
  interfaces y decisiones técnicas a partir de la especificación aprobada.
  \emph{Resultado:} diseños revisados y relacionados con sus requisitos.

  \item \textbf{Implementación e integración.} Construir los componentes e
  integrarlos respetando las interfaces acordadas. \emph{Resultado:}
  versión integrada del sistema.

  \item \textbf{Verificación, validación y aceptación.} Ejecutar pruebas
  de componentes, integración, sistema y aceptación; contrastar los
  resultados con requisitos y criterios previamente definidos.
  \emph{Resultado:} evidencias de conformidad y registro de defectos.

  \item \textbf{Despliegue y mantenimiento.} Instalar, operar, corregir
  defectos y tramitar nuevas necesidades mediante solicitudes de cambio
  evaluadas por su impacto. \emph{Resultado:} sistema operativo con
  documentación y versiones controladas.
\end{enumerate}

Las fases describen una secuencia predominante, no la prohibición absoluta
de retroceder: una no conformidad puede exigir revisar una fase anterior,
pero el cambio debe evaluarse y documentarse.

\subsection{Supuestos del modelo}
Para que esta organización sea razonable, se parte de los siguientes
supuestos, que deben comprobarse en el proyecto concreto:
\begin{enumerate}
  \item Las necesidades fundamentales pueden conocerse y aclararse antes
  de comprometer el diseño detallado.
  \item Los requisitos centrales y las interfaces tendrán una estabilidad
  relativa durante la construcción.
  \item Los responsables pueden revisar y aprobar entregables en hitos
  definidos, aunque no participen todos los días.
  \item Es viable especificar criterios de aceptación y conservar
  evidencias de cumplimiento para cada requisito relevante.
  \item Los cambios posteriores se gestionarán mediante un procedimiento
  explícito de análisis de impacto, aprobación y actualización documental.
\end{enumerate}

\subsection{Condiciones de aplicación y límites}
\textbf{Es favorable} cuando el alcance está delimitado, existen interfaces
técnicas conocidas, hay exigencias de documentación y trazabilidad, las
revisiones formales son importantes y el cliente puede aprobar una línea
base antes del desarrollo \cite{swebok2024}.

\textbf{Es poco conveniente como proceso estrictamente secuencial} cuando
los objetivos cambian con frecuencia, se desconoce la viabilidad técnica,
se necesita experimentación continua o las personas usuarias solo pueden
aclarar sus expectativas al probar prototipos. En estos casos, conviene
considerar un enfoque iterativo o introducir una exploración técnica
previa a la aprobación de la línea base.


